\section{Auld Lang Syne}

\begin{quotex}
To live without principles is to be unprincipled.

To act without a point is to be pointless.
\end{quotex}


The end of the year is customarily the time to review the past year and plan for the next. Because of the personal equation, there is no yellow brick road to follow. Instead, a man must be a wanderer, following instincts and intuitions.

During my end of year office cleanup, I found some personal journals, from 1978 through 1982. There was a collection of speeches from 1992. In particular, one was on climate change which could have been presented last weekend without revision; particular its relationship to Santa Claus. Another was based on certain ideas from Bernard Lonergan; I'll edit it a little and put it on line.

The topics were the same as today for the most part. I wrote out goals, study plans, and outlines for essays. In those days I was more nihilistic. Probably still am because ultimately there is no unassailable proof for anything. Heck, I don't even know how many, if any, parallel lines there are that go through a point not on another line.

You have to choose; the proof is in the doing. That is because Knowing is Being. You can't know unless you are. Thomas Aquinas writes:

\begin{quotex}
Nothing is known except Truth, which is the same as Being.
\end{quotex}


Existence is the actualization of the possibilities of Being. The actualization of the soul's potentiality is to become what it contemplated. Contemplate God, and become God: this is \emph{theosis}.

\paragraph{Influences}


There was a short section on five people or ideas who were influencing me.

\begin{itemize}
\item \textbf{William Blake}: \emph{Those who restrain their desire do so only because it is weak enough to restrain.}

This is essential to Western thought (e.g., drive is a major element in Spinoza). Perhaps I have confused the Buddhist idea of the elimination of desire with the repression of desire, but I have reached the realization that inhibition is contrary to full self-expression.

\emph{Comment 2020}: This is still true. Those who are too weak to act then call it virtue.

\item \textbf{Weekend Seminar}: The notion of Expansion vs Contraction. I have counting on my Good to come to me (a contracting idea). I now work on myself to expand and encompass more of my own good.

\emph{Comment 2020}: I still stand by this. A popular technique is visualization or affirmations whose practitioners beg the ``universe'' to satisfy their needs and desires. That is a contracting approach.

That is opposite to the way the universe ``works''; since the Fall, our needs are met through toil and suffering. The outcome is uncertain because of death. Everything we create is subject to death and decay; hence, a new spark becomes necessary. You expand or disappear.

\item \textbf{Weekend Seminar 2}: the idea that we choose our own experience.

\emph{Comment 2020}: This may sound trite, but assume it is true as the null hypothesis. Think carefully and deeply about your life. See how you were implicated in much of what happened to you. Obviously, not everything since there are indeed accidents. Nevertheless, probably more than you suppose.

\item \textbf{Mircea Eliade}: on the world as the Absolute (in Parabola)

\emph{Comment 2020}: I used to have a full collection of Parabola magazines up through the 1990s. I don't recall this particular article and the Parabola archive is not helpful. I'm assuming it came from \emph{The Sacred and the Profane}. Here are some quotes.

This is as much as to say that every religious man places himself at the Center of the World and by the same token at the very source of absolute reality, as close as possible to the opening that ensures him communication with the gods

Whatever the historical context in which he is placed, \emph{homo religiosus} always believes that there is an absolute reality, the sacred, which transcends this world but manifests itself in this world, thereby sanctifying it and making it real.

Perceived by virtue of a religious experience, the specific mode of existence of the stone reveals to man the nature of an absolute existence, beyond time, invulnerable to be. coming.

\item \textbf{My Father}: Taught me to develop the strength not to succumb to fear and that any decision made, for better or for worse, must be followed through, despite anxiety and public opinion.

\emph{Comment 2020}. This is particularly helpful to me whenever unresolved Oedipal issues arise. I still live by it, particularly the absoluteness of decisions.
\end{itemize}

\paragraph{Knights}

Knights venture out to seek adventure and impress their ladies fair. Nothing essential has changed. The adventures are different, but so are the women.

The original knights were free. A free man has the right to bear arms and to spend his money the way he likes. The knight would entertain his friends and build alliances.

The old ways were built on Loyalty, Omertà, Respect, Honor.

Only the free man can pledge loyalty. The loyalty of the serf is compelled; he trades freedom for protection and security.

Few understand the art of gift giving today. When invited to dinner, I bring a gift. Perhaps PBR and a box of donuts, or else a bottle of wine and French pastries, depending on the host. I don't hold onto the gift until after the meal, and then give it to the hostess only if I enjoyed the meal.

It's best to give a gift that the recipient could get in no other way. A gift carries no expectations.

A knight was chivalrous. Like everything noble, when a task filters down to the bourgeoisie, it deteriorates. Woman today don't see the value in it, other than pro forma acts.

One woman never tired of repeating that if her ex-husband ever needed her, she would rush to his aid. She even repeated her pledge to a salesman and asked if he would like an ex-wife like that, and he agreed. I then asked if he wanted a gf who would return to her husband at a moment's notice. Of course not, so from that point, I treated her like a friend with benefits. There was no point to argue about it.

One woman asked me for a song to describe our relationship. I sent her this one: I will follow him\footnote{\url{https://www.youtube.com/watch?v=jgPMYQTINNk}}. She scoffed at me, as if she had a long queue of better suitors. She was unaware of the benefits such loyalty would bestow on her. So I ghosted her.

Knights have no mercy to those who don't follow the code.

\paragraph{Experiencing the World}


\begin{quotex}
Transcending oneself: this is the great imperative of the human condition; and there is another that anticipates it and at the same time prolongs it: dominating oneself. The noble man is the one who dominates himself; the holy man is the one who transcends himself.
\flright{\textsc{Frithjof Schuon}}
\end{quotex}


A benefit of living like an anchorite is that I been able to devote 1 to 2 hours per day for Meditation time. You develop a sense of detachment of the Self. Live is not lived but is observed instead. The Self watches the body squirm around a bit until it finds a comfortable position. You watch is fall asleep; sometimes you can watch it wake up.

The soul life is like a rippling pond, each ripple being a perturbation of consciousness. A feeling arises, a thought appears, a fantasy presents itself. Sometimes the ripples become waves and they need to be calmed.

All your emotions, all hidden motives are exposed; it can be difficult to act on principles and to avoid pointless behavior. But you are seldom plagued with feelings of dread, anxiety, fear of death.

The light of the Holy Spirit is best reflected on still water.

\paragraph{The Polemic Spirit}


Rational discourse is an inferior form of knowledge. Nevertheless, a certain type of fellow thinks the secret to enlightenment lies in undirected study and disputations. He lives in his head, like a brain in a vat. The Germans have a colourful word for that: Kopfmensch.

Avoid those with a polemic spirit. The esoteric path is one of harmonizing disparate points of view.

Rene Guenon, in a private letter, reveals more of himself than he does in his books:

\begin{quotex}
Meditation is more important than reading, which can only serve as a starting point, and it is generally beneficial not to overdo it in order to avoid dissipation. Rites are effective in themselves, but of course, attention and concentration strengthen the effects significantly. I particularly recommend the regular recitation of the \emph{wird} Sufi liturgy, because this especially strengthens the link with the tarîqa. Finally, I think everyone should seek to use their natural tendencies rather than fight them; but, of course, there are as many different modalities as there are individualities.
\flright{\textit{Letter from Rene Guenon to L.C. d'Amiens}, 30 August 1935}
\end{quotex}


Metaphysical doctrine can be taught through the intellect, but it can only be understood through concentration, meditation, and contemplation. There is no one universal path to metaphysical realization. Hafiz, for example, did a vigil at a saint's tomb for 40 nights. Today, there are many who think they desire such lofty realization, yet don't have one hour per week to spare.

The point is that in order to describe a path, I need to reveal personal details. Many are not flattering, but I hope that are helpful. A wanderer will encounter many things and do many things. Perhaps, like an alchemist, he can transform dross into gold. Or die trying.

\paragraph{Un bel Giorno per Morire}


I asked my nephrologist about the process of dying from kidney disease: was it quick or slow and painful? He said the latter but reassured me that his goal is to make sure I die of something other than kidney disease. The following week I met with my cardiologist. She explained that her goal is to make sure that I die of something other than heart disease. I think the three of us need to get together and work out the details.

Sometimes I feel nostalgic and wonder what I'll miss most about the world. Young children delight in every new experience. They demand, ``again, again''. When is enough, enough, and repetition loses its flavour? So much of what once seemed important really counts for nothing. The World Process will continue.

Repetition has led to the forgetfulness of the awesome significance of forgiveness. It counters free will and stops human karma; otherwise, things get worse and worse. God the Merciful forgets our forgiven sins; our unforgiven sins are judged harshly.

It is a beautiful day to die\footnote{\url{https://www.youtube.com/watch?v=RnCmeNa5vXs}} and in a leap year we have 366 such days.

There are so many good versions, quite different from each other, of Robert Burn's poem, it was difficult to choose one.

\flrightit{Posted on 2020-12-31 by Cologero}

